\documentclass[11pt,a4paper]{article}
\usepackage[UTF8]{ctex}
\usepackage{geometry}
\geometry{left=2.5cm,right=2.5cm,top=2.5cm,bottom=2.5cm}
\usepackage{amsmath,amssymb,amsthm}
\usepackage{hyperref}
\usepackage{booktabs}
\usepackage{enumitem}
\usepackage{xltabular}
\hypersetup{colorlinks=true,linkcolor=blue,urlcolor=blue}

\title{基于生成论的网络安全分析能力调度系统\\
附录}
\author{李龙胜}
\date{\today}

\begin{document}

\maketitle

\begin{center}
\textit{
符号的锚点、公式的源头、常见问题的回答、许可证的全文。\\
这是整个理论体系的工具抽屉。
}
\end{center}

\vspace{5mm}

\tableofcontents

\section{附录A：符号索引}

本索引覆盖全四卷使用的关键符号。每个符号的正式锚点见维护工具包。

\subsection{告警分级与分析能力}

\begin{center}
\begin{xltabular}{\textwidth}{c|X}
\toprule
\textbf{符号} & \textbf{含义} \\
\midrule
$r_i$ & 告警类别$i$的分析率（被分析的概率），$r_i \in [0,1]$ \\
$r^*$ & 最优分析率，由贪心择优裁定 \\
$A_i$ & 告警类别$i$攻击成功时的最大损失（货币化度量） \\
$\alpha_i$ & 告警类别$i$对漏报的敏感度系数 \\
$\beta_i$ & 告警类别$i$的结构特征常数，$\beta_i = \alpha_i A_i^2 / (2B_i)^2$ \\
$B_i$ & 告警类别$i$的攻击手法变化率（带宽） \\
$\kappa$ & 单条告警的单位分析成本（时薪$\times$耗时） \\
$C$ & 每日总分析能力预算（容量上限） \\
$C_{\text{consumed}}$ & 已消耗的分析能力（实部） \\
$C_{\text{reserve}}$ & 保留的分析能力（虚部） \\
$\mathcal{K}_{\text{covered}}$ & 当前分析能力有效覆盖的告警类型集合 \\
$V_{\text{breadth}}$ & 虚部广度，$|\mathcal{K}_{\text{covered}}| / K$ \\
$V_{\text{risk}}$ & 虚部风险暴露，未被覆盖类型的潜在损失之和 \\
$q_{\min}$ & 同源关联优先策略的生效阈值 \\
\bottomrule
\end{xltabular}
\end{center}

\subsection{柯克霍夫原则与安全定义}

\begin{center}
\begin{xltabular}{\textwidth}{c|X}
\toprule
\textbf{符号} & \textbf{含义} \\
\midrule
$\theta$ & 防御方的秘密参数集合（$A_i^{\min}, \alpha_i$等） \\
$\varepsilon$ & 参数不可区分性的安全参数（越小越安全） \\
$\delta$ & 参数周期性扰动的幅度 \\
$T_{\text{rot}}$ & 参数扰动周期 \\
$q_{01}$ & 蓝队伪造概率：将忽略伪装为分析 \\
$q_{10}$ & 蓝队压制概率：将分析伪装为忽略 \\
$\Delta$ & 两个参数集合$\theta_0$和$\theta_1$对应的可观测差异 \\
$M_{\max}$ & 攻击方对单个告警类型的探测预算上限 \\
$m^*$ & 区分参数所需的最小探测次数（信息论下界） \\
\bottomrule
\end{xltabular}
\end{center}

\subsection{攻击方与博弈}

\begin{center}
\begin{xltabular}{\textwidth}{c|X}
\toprule
\textbf{符号} & \textbf{含义} \\
\midrule
$\ell$ & 攻击方的质量等级，$\ell \in [0,1]$（自动化程度与隐蔽性的权衡） \\
$C_A$ & 攻击方总资源预算（算力 + IP池） \\
$q_{\text{hit}}$ & 单次攻击尝试命中防御方脆弱点的概率 \\
$\rho$ & 攻击方对信息暴露的敏感度系数 \\
$s(t)$ & 防御方分析能力饱和度，$s(t) = C_{\text{consumed}}(t)/C$ \\
$\tau(t)$ & 防御方响应延迟（饱和指标之一） \\
$\theta_{mn}$ & 攻击方$m$和$n$之间的协同参数，$\theta_{mn} \in [0,1]$ \\
\bottomrule
\end{xltabular}
\end{center}

\subsection{生成步系统（数学基础）}

\begin{center}
\begin{xltabular}{\textwidth}{c|X}
\toprule
\textbf{符号} & \textbf{含义} \\
\midrule
$R_{\min}$ & 最小记录量（公设二的数学表达） \\
$s_T$ & 系统在逻辑时间$T$的状态 \\
$U$ & 生成步（作用在状态上的操作） \\
$R(s,U)$ & 记录成本函数 \\
$\mathcal{A}(s)$ & 状态$s$处允许的生成步集合 \\
$J(\pi)$ & 路径$\pi$的累积记录成本 \\
$f(s)$ & 单调进展函数（贪心最优性定理的条件） \\
\bottomrule
\end{xltabular}
\end{center}

\section{附录B：数学推导补遗}

\subsection{B.1 最优采样频率方程的推导（Math-51 核心公式）}

设信号 $x(t)$ 有效带宽为 $B$，动态范围为 $A$，采样频率为 $f_s$，奈奎斯特频率为 $2B$。
定义无量纲分析率 $r \equiv f_s / (2B)$。

当 $r < 1$ 时，混叠误差导致的状态偏差成本为：
\[
V(r) = \alpha A^2 \left( \frac{1}{r} - 1 \right)^2, \quad (r < 1).
\]
评估成本为 $K(r) = \kappa f_s = 2B\kappa r$。

总成本 $J(r) = V(r) + K(r)$。对 $r$ 求导并令其为零：
\[
\frac{dJ}{dr} = -2\alpha A^2 \left( \frac{1}{r} - 1 \right) \frac{1}{r^2} + 2B\kappa = 0.
\]
整理得：
\[
\frac{\alpha A^2}{r^3} \left( \frac{1}{r} - 1 \right) = B\kappa.
\]
引入信号特征常数 $\beta \equiv \frac{\alpha A^2}{(2B)^2}$，替换后得到无量纲方程：
\[
\frac{1-r}{r^4} = \frac{\kappa}{2\beta}.
\]
该方程在 $(0,1]$ 内有唯一解 $r^*$。最优采样频率为 $f_s^* = 2B \cdot r^*$。

\subsection{B.2 贪心最优性定理的完整证明骨架}

\textbf{定理}（贪心最优性定理）：
设生成步系统 $\mathcal{G}$ 满足：
\begin{enumerate}[label=(C\arabic*)]
    \item 存在严格单调进展函数 $f: \mathcal{S} \to \mathbb{R}$，每步进展至少为 $\delta > 0$。
    \item 成本函数可分解为 $R(s,U) = R_{\text{base}}(f(s)) + R_{\text{trans}}(\Delta f)$，
          其中 $R_{\text{base}}$ 单调不减，$R_{\text{trans}}$ 严格凸且 $R_{\text{trans}}(0)=0$。
    \item 累积记录成本不触及容量上限 $C$。
\end{enumerate}
则对任何初始状态 $s_0$，局部贪心路径也是全局累积成本最小的路径。

\textbf{证明}（交换论证骨架）：

设贪心路径为 $\pi^* = (U_0^*, U_1^*, \ldots)$，假设存在另一条路径 $\pi$ 使 $J(\pi) < J(\pi^*)$。
令 $T$ 为两条路径首次分歧的时刻。在时刻 $T$，贪心选择的进展量为 $\Delta_T^*$，
$\pi$ 选择的进展量为 $\Delta_T$，由贪心定义有 $R(s_T, U_T^*) \le R(s_T, U_T)$。

构造新路径 $\tilde{\pi}$：在第 $T$ 步取 $U_T^*$，从 $T+1$ 步开始，模仿 $\pi$ 的后续步，
但将每步的进展量按比例调整以保持总进展量一致。由 $R_{\text{trans}}$ 的严格凸性，
$\tilde{\pi}$ 的总成本严格小于 $\pi$ 的总成本。重复此交换直至消除所有分歧，
可得 $J(\pi^*) \le J(\pi)$，矛盾。故贪心路径最优。

\subsection{B.3 信息论下界的详细推导}

设攻击方对告警类型 $i$ 进行 $m$ 次独立探测，每次探测伯努利试验，
真实分析概率为 $P_{\text{obs}}$。攻击方用频率 $\hat{P} = X/m$ 估计 $P_{\text{obs}}$。

由 Chernoff-Hoeffding 不等式，对任意 $\varepsilon > 0$：
\[
\Pr\left(|\hat{P} - P_{\text{obs}}| \ge \varepsilon\right) \le 2e^{-2m\varepsilon^2}.
\]
要使得估计误差超过 $\varepsilon$ 的概率不超过 $\delta$，即 $2e^{-2m\varepsilon^2} \le \delta$，
解得：
\[
m \ge \frac{1}{2\varepsilon^2} \ln\frac{2}{\delta}.
\]
代入 $\varepsilon = 0.01, \delta = 0.001$，得 $m \ge 38000$。

攻击方要区分 $\theta_0$ 和 $\theta_1$，需满足 $\Delta = |P_{\text{obs}}(\theta_0) - P_{\text{obs}}(\theta_1)|$。
区分优势满足 $\Pr[\text{猜对}] \approx \Phi(\Delta\sqrt{m})$，其中 $\Phi$ 是标准正态累积分布。
要达 90\% 成功率需 $\Delta\sqrt{m} \ge 1.28$，即 $m \ge (1.28/\Delta)^2$。
当 $\Delta = 0.01$ 时 $m \ge 16384$，与 Chernoff 下界一致。

\section{附录C：常见问题（FAQ）}

\subsection*{C.1 这和设备的威胁评级有什么区别？}
设备评的是“这个攻击有多危险”，我们算的是“这个告警值不值得花时间去查”。
一个高危告警如果每次都自动被阻断，查它就是浪费人力。

\subsection*{C.2 如果我漏掉了一条告警导致出事，谁负责？}
这套算法的逻辑是帮您把时间花在最可能出事的地方。
如果出事，出事的告警一定是被算法标注为“低分析率”的——
那您需要反思的是：为什么这个本该是高分析率的告警，您当初给它标了低损失值？

\subsection*{C.3 需要什么样的技术栈？}
只需要三种东西：一张台账表（Excel或MySQL）、一个可以运行Python脚本的服务器、
以及安全设备的数据接口。不需要机器学习，不需要大数据平台。

\subsection*{C.4 我的团队很小，这样做有意义吗？}
团队越小，越需要把力气用在刀刃上。小团队根本不可能处理所有告警。

\subsection*{C.5 初始参数怎么填？}
初期手工分档：核心服务器告警 $A=10^6$ 元，普通终端 $A=10^4$ 元，
自动阻断告警 $A=10^2$ 元。$\alpha$ 统一设为 1.0。
后续依靠自适应反馈机制自动校准。

\subsection*{C.6 柯克霍夫原则要求参数保密，但合规审计怎么办？}
参数本身加密存储，审计时可以证明确实存在参数管理和扰动机制，
但不需要公开参数的精确值。与密码学中密钥管理类似：
可以证明密钥存在且被妥善管理，而不需要公开密钥本身。

\section{附录D：CC BY 4.0 许可证全文}

本项目全部文档（包括全四卷及本附录）采用知识共享署名 4.0 国际公共许可证（CC BY 4.0）。

\subsection*{您可以自由地：}
\begin{itemize}
    \item \textbf{共享} — 在任何媒介以任何形式复制、发行本作品。
    \item \textbf{改编} — 修改、转换或以本作品为基础进行创作。
    \item 用于任何目的，包括商业目的。
\end{itemize}

\subsection*{惟须遵守下列条件：}
\begin{itemize}
    \item \textbf{署名} — 您必须给出适当的署名，提供指向本许可协议的链接，
          同时标明是否对原始作品作了修改。您可以用任何合理的方式来署名，
          但不得以任何方式暗示许可人为您或您的使用背书。
\end{itemize}

\textbf{CC BY 4.0 完整法律文本参见：}\\
\url{https://creativecommons.org/licenses/by/4.0/legalcode}

\vspace{5mm}
\noindent\textbf{文档信息}：
\begin{itemize}
    \item 作者：李龙胜
    \item 版本：v1.0（与四卷体系同步发布）
    \item 许可：CC BY 4.0
\end{itemize}

\end{document}